%% ============================================================
%%  INRC-II Contextual Bandits — Final Report
%%  DS 592: Introduction to Sequential Decision Making
%%  Spring 2026
%%  Team: Ria Sonalker · Jihyeon Yun · Achuthan Rathinam
%% ============================================================

\documentclass[12pt, a4paper]{article}

% ── Encoding & fonts ─────────────────────────────────────────
\usepackage[T1]{fontenc}
\usepackage[utf8]{inputenc}
\usepackage{lmodern}
\usepackage{microtype}

% ── Page layout ──────────────────────────────────────────────
\usepackage[top=1in, bottom=1in, left=1.25in, right=1.25in]{geometry}
\usepackage{setspace}
\onehalfspacing

% ── Math ─────────────────────────────────────────────────────
\usepackage{amsmath, amssymb, amsthm}
\usepackage{mathtools}
\usepackage{bm}          % bold math symbols

% ── Algorithms ───────────────────────────────────────────────
\usepackage[ruled, vlined, linesnumbered]{algorithm2e}

% ── Tables & figures ─────────────────────────────────────────
\usepackage{booktabs}
\usepackage{tabularx}
\usepackage{multirow}
\usepackage{graphicx}
\usepackage{subcaption}
\usepackage{float}

% ── Listings / code ──────────────────────────────────────────
\usepackage{listings}
\usepackage{xcolor}
\lstset{
  basicstyle=\ttfamily\small,
  keywordstyle=\color{blue},
  commentstyle=\color{gray},
  stringstyle=\color{teal},
  breaklines=true,
  frame=single,
  numbers=left,
  numberstyle=\tiny\color{gray},
}

% ── References & links ───────────────────────────────────────
\usepackage[hidelinks]{hyperref}
\usepackage{cleveref}
\usepackage[numbers, sort&compress]{natbib}

% ── Misc ─────────────────────────────────────────────────────
\usepackage{enumitem}
\usepackage{parskip}
\usepackage{lipsum}      % remove when real content is in place

% ── Custom commands ──────────────────────────────────────────
\newcommand{\xt}{\mathbf{x}_t}          % context vector at time t
\newcommand{\arm}[1]{A_{#1}}             % arm shorthand
\newcommand{\reward}{r_t}               % reward at time t
\newcommand{\policy}{\pi}               % policy
\newcommand{\regret}{\mathcal{R}_T}     % cumulative regret up to T
\newcommand{\TODO}[1]{\textbf{\textcolor{red}{[TODO: #1]}}}

% ── Theorem environments ─────────────────────────────────────
\newtheorem{theorem}{Theorem}[section]
\newtheorem{lemma}[theorem]{Lemma}
\newtheorem{definition}{Definition}[section]
\newtheorem{remark}{Remark}[section]

% ============================================================
\title{
  \vspace{-1cm}
  \Large\textbf{Nurse Rostering as a Contextual Bandit Problem:\\
  A Reimplementation Study on INRC-II}\\[0.5em]
  \normalsize DS 592: Introduction to Sequential Decision Making\\
  Spring 2026
}
\author{
  Ria Sonalker \and Jihyeon Yun \and Achuthan Rathinam
}
\date{\today}

% ============================================================
\begin{document}

\maketitle
\thispagestyle{empty}

% ── Abstract ─────────────────────────────────────────────────
\begin{abstract}
\noindent
\TODO{Write abstract once reimplementation is complete.}
%
% Template: This report presents our reimplementation of [PAPER TITLE] on the
% Second International Nurse Rostering Competition (INRC-II) benchmark.
% We frame the rostering task as a contextual bandit problem: at each week t,
% a context vector summarising nurse history and coverage demand is observed,
% one of K scheduling heuristics (arms) is selected, and a reward based on
% soft-constraint penalties is received.  We reproduce the key empirical
% findings of the original paper and discuss divergences, ablations, and
% insights gained from our implementation.
\end{abstract}

\newpage
\tableofcontents
\newpage

% ============================================================
\section{Introduction}
\label{sec:intro}
% ── Purpose of this section ──────────────────────────────────
% • Motivate the problem: why nurse rostering, why bandits?
% • State the paper being reimplemented and what claims you target.
% • Brief outline of the report structure.
% ─────────────────────────────────────────────────────────────

Nurse rostering is a combinatorial optimisation problem in which a set of
nurses must be assigned to shifts over a planning horizon while respecting
hard contractual constraints and minimising soft-preference violations.
The Second International Nurse Rostering Competition (INRC-II)~\cite{inrc2}
provides a standardised benchmark with realistic multi-week instances.

This report documents our reimplementation of \TODO{[paper citation here]}.
We treat each week of the planning horizon as a round of a contextual bandit
game: the learner observes a context $\xt \in \mathbb{R}^d$ encoding the
current workforce state and coverage demand, selects one of $K$ scheduling
heuristics (arms), and receives an immediate reward $\reward$ reflecting
constraint satisfaction.

The remainder of this report is organised as follows:
\Cref{sec:background} reviews the INRC-II problem and relevant bandit
algorithms.
\Cref{sec:paper} summarises the paper being reimplemented.
\Cref{sec:formulation} gives our formal problem formulation.
\Cref{sec:method} describes the reimplemented methodology.
\Cref{sec:experiments} presents experimental results.
\Cref{sec:discussion} discusses findings and deviations.
\Cref{sec:conclusion} concludes.

% ============================================================
\section{Background}
\label{sec:background}

\subsection{INRC-II Problem Description}
% • Scenario / week / history file structure.
% • Hard constraints (forbidden successions, coverage minimums).
% • Soft constraints and penalty weights.

\TODO{Summarise INRC-II problem — refer to \texttt{ref\_docs/INRC2.pdf}.}

\subsection{Contextual Bandits}
% • Define the contextual bandit model (context, action, reward).
% • Cover relevant algorithms: LinUCB, Thompson Sampling, Neural Bandits, etc.
% • State regret bounds if relevant to the paper.

\begin{definition}[Contextual Bandit]
A contextual bandit is a tuple $(\mathcal{X}, \mathcal{A}, r)$ where
$\mathcal{X}$ is a context space, $\mathcal{A} = \{1, \ldots, K\}$ is a
finite action set, and $r : \mathcal{X} \times \mathcal{A} \to \mathbb{R}$
is a (stochastic) reward function.  At each round $t$ the learner observes
$\xt \in \mathcal{X}$, selects action $a_t \in \mathcal{A}$ according to
policy $\policy$, and receives $\reward = r(\xt, a_t) + \epsilon_t$.
\end{definition}

\TODO{Expand with algorithm summaries once the target paper is confirmed.}

% ============================================================
\section{Paper Being Reimplemented}
\label{sec:paper}
% ── Fill this section once you feed the paper ────────────────
% • Full citation.
% • Core claims / contributions.
% • Algorithm(s) proposed.
% • Evaluation protocol and reported results.
% ─────────────────────────────────────────────────────────────

\TODO{This section will be filled after the target paper is provided.}

\subsection{Core Claims}

\subsection{Proposed Algorithm}

\subsection{Evaluation Protocol and Reported Baselines}

% ============================================================
\section{Problem Formulation}
\label{sec:formulation}

\subsection{State / Context}

At the start of week $t \in \{1, \ldots, W\}$ we construct a context vector
\begin{equation}
  \xt = \bigl[\mathbf{x}^{\text{nurse}}_t \;\|\; \mathbf{x}^{\text{week}}_t\bigr]
  \in [0,1]^d
\end{equation}
where the nurse-level features aggregate per-nurse contract utilisation
(consecutive days, weekends worked, total assignments) and the week-level
features capture coverage demand, day-off requests, and horizon position.

\subsection{Arms}

We define $K = 5$ scheduling heuristics as arms:

\begin{itemize}[leftmargin=2em]
  \item $\arm{1}$: \textbf{Coverage-First Greedy} — fill largest coverage gaps first.
  \item $\arm{2}$: \textbf{Fatigue-Aware} — prioritise rest for nurses near consecutive-day limits.
  \item $\arm{3}$: \textbf{Weekend-Balancing} — minimise variance in weekends worked.
  \item $\arm{4}$: \textbf{Preference-Respecting} — maximise satisfied day-off requests.
  \item $\arm{5}$: \textbf{Lookahead-Conservative} — prefer nurses with the most remaining contract slack.
\end{itemize}

\subsection{Reward Signal}

The per-week proxy reward is defined as
\begin{equation}
  \reward = -\bigl(\lambda_1 p_{\text{cov}} + \lambda_2 p_{\text{pref}}
            + \lambda_3 p_{\text{consec}} + \lambda_4 p_{\text{succ}}\bigr)
\end{equation}
where $p_{\text{cov}}$, $p_{\text{pref}}$, $p_{\text{consec}}$,
$p_{\text{succ}}$ are soft-constraint penalty counts and
$\lambda_i$ are tunable weights.

At the final week $W$, the official Java validator is run and a correction
term $\delta$ is applied: $r_W \mathrel{+}= \delta$.

\subsection{Objective}

Minimise cumulative regret over $T = N \cdot W$ rounds (where $N$ is the
number of training instances):
\begin{equation}
  \regret = \sum_{t=1}^{T} \bigl[r^*_t - r_t\bigr]
\end{equation}
where $r^*_t = \max_{a \in \mathcal{A}} \mathbb{E}[r(\xt, a)]$ is the
expected reward of the optimal arm given context $\xt$.

% ============================================================
\section{Methodology}
\label{sec:method}

\subsection{Input Parsing}
% Scenario XML, Week XML, History JSON → Python dataclasses

\TODO{Describe parser design.}

\subsection{Feature Engineering}
% Normalisation, bound fitting on training set only

\TODO{Describe feature engineering pipeline.}

\subsection{Bandit Algorithm}
% Which algorithm, hyperparameters, training loop

\TODO{Fill once target paper algorithm is confirmed.}

\begin{algorithm}[H]
\SetAlgoLined
\KwIn{Instances $\mathcal{I}$, arms $\mathcal{A}$, horizon $W$}
\KwOut{Learned policy $\hat{\policy}$}
Initialise bandit model\;
\For{each instance $i \in \mathcal{I}$}{
  \For{$t = 1$ \KwTo $W$}{
    Observe context $\xt$\;
    Select arm $a_t \leftarrow \policy(\xt)$\;
    Execute heuristic $\arm{a_t}$, obtain roster\;
    Compute reward $\reward$\;
    Update bandit model with $(\xt, a_t, \reward)$\;
  }
}
\caption{Contextual Bandit for INRC-II Rostering}
\label{alg:bandit}
\end{algorithm}

\subsection{Feasibility Repair}
% Forbidden successions, minimum coverage, greedy swap

\TODO{Describe repair heuristic.}

\subsection{Evaluation Protocol}
% Train / test split, instance sets, evaluation metric

\TODO{Describe evaluation protocol.}

% ============================================================
\section{Experiments}
\label{sec:experiments}

\subsection{Setup}

\TODO{Hardware, software versions, random seeds.}

\subsection{Baseline Comparison}

\begin{table}[H]
  \centering
  \caption{Cumulative soft-constraint penalty (lower is better) on test instances.
           \TODO{Fill with real numbers.}}
  \label{tab:results}
  \begin{tabularx}{\linewidth}{lXXX}
    \toprule
    Method & Instance Set A & Instance Set B & Instance Set C \\
    \midrule
    Random arm & -- & -- & -- \\
    Best fixed arm (oracle) & -- & -- & -- \\
    \TODO{Paper method} & -- & -- & -- \\
    \textbf{Our reimpl.} & -- & -- & -- \\
    \bottomrule
  \end{tabularx}
\end{table}

\subsection{Ablations}

\TODO{Feature ablation, reward weight sensitivity, arm subset experiments.}

\subsection{Learning Curves}

\TODO{Plot cumulative regret vs.\ rounds.}

% ============================================================
\section{Discussion}
\label{sec:discussion}

\subsection{Reproduction of Original Results}

\TODO{State which claims were reproduced and which deviated. Explain divergences.}

\subsection{Implementation Challenges}

\subsection{Limitations}

% ============================================================
\section{Conclusion}
\label{sec:conclusion}

\TODO{Summarise findings, lessons learnt, and future directions.}

% ============================================================
\appendix

\section{Heuristic Pseudocode}
\label{app:heuristics}

\TODO{Pseudocode for each of the 5 arms.}

\section{Hyperparameter Settings}
\label{app:hyperparams}

\begin{table}[H]
  \centering
  \caption{Hyperparameter settings used in all experiments.}
  \begin{tabular}{ll}
    \toprule
    Parameter & Value \\
    \midrule
    \TODO{param} & \TODO{value} \\
    \bottomrule
  \end{tabular}
\end{table}

% ============================================================
\bibliographystyle{plainnat}
\bibliography{references}

\end{document}
